\chapter*{ВВЕДЕНИЕ}
\addcontentsline{toc}{chapter}{ВВЕДЕНИЕ}

При отборе фотографий для экспертного анализа специалисту нередко доступна только последняя JPEG-копия. Снимок могли сохранить при экспорте из редактора или при подготовке к публикации. Внешне две версии могут выглядеть одинаково, хотя второе сжатие изменило распределение частотных коэффициентов. По этим изменениям можно выделить часть повторно сохранённых файлов и передать их специалисту для дальнейшего анализа~\cite{wang2016,park2018}.

Цель работы --- разработать программу, которая обнаруживает признаки повторного JPEG-сжатия по одному файлу, и установить, как выбор нейросетевой модели влияет на качество классификации и время обработки изображения.

Для достижения цели поставлены следующие задачи.
\begin{enumerate}
  \item Сформулировать требования к программе и ограничения на входные изображения.
  \item Сравнить методы обнаружения повторного JPEG-сжатия и выбрать нейросетевые модели для исследования.
  \item Спроектировать структуру программы и алгоритмы подготовки данных, обучения и классификации.
  \item Выбрать средства реализации и разработать программу для обработки JPEG-файлов.
  \item Обучить модели на размеченном наборе однократно и двукратно сжатых изображений.
  \item Сопоставить ошибки и время обработки у двух моделей, определить влияние параметров первого и второго сжатия на результат.
\end{enumerate}
